\documentclass[11pt]{article}

\usepackage[margin=1in]{geometry}
\usepackage{amsmath,amssymb,amsthm,mathtools}
\usepackage{hyperref}

\newcommand{\R}{\mathbb{R}}
\newcommand{\tr}{\operatorname{tr}}

\newtheorem{theorem}{Theorem}
\newtheorem{proposition}[theorem]{Proposition}
\newtheorem{corollary}[theorem]{Corollary}
\theoremstyle{remark}
\newtheorem{remark}[theorem]{Remark}

\title{Why Initialization Alone Is Not Enough: Stationary-Law Design for Finite-\texorpdfstring{$K$}{K} BatchEnsemble}
\author{Research note generated in sandbox}
\date{2026-04-05}

\begin{document}
\maketitle

\section{Message}

The local contrast-mode theory gives principled initial designs. However, the user's objection is exactly right:

\begin{quote}
why should an initialization that is locally optimal near the start remain optimal after long nonlinear training?
\end{quote}

In general, it should not. To justify a finite-$K$ design for the final predictor, one must reason about the \emph{stationary law} or terminal training state, not just the initialization. This note makes that statement precise and gives a two-stage design principle based on the converged collapsed law.

\section{Initialization-Only Design Can Be Washed Out}

Consider nonlinear member dynamics of the head-local form
\begin{equation}
\dot q_{\alpha,t} = \mathcal C(\rho_t,q_{\alpha,t}),
\qquad
\alpha=1,\dots,K,
\label{eq:member-dynamics}
\end{equation}
where $\rho_t$ denotes the shared state. We make no linearization here.

\begin{theorem}[Nonlinear contraction destroys initialization-only diversity]
\label{thm:nonlinear-contraction}
Assume there exists $\gamma>0$ such that for every shared state $\rho$ and every pair $q,q'\in\R^p$,
\begin{equation}
\langle q-q',\ \mathcal C(\rho,q)-\mathcal C(\rho,q')\rangle
\le
-\gamma \|q-q'\|^2.
\label{eq:one-sided-lipschitz}
\end{equation}
Then for every pair of members $\alpha,\beta$,
\[
\|q_{\alpha,t}-q_{\beta,t}\|
\le
e^{-\gamma t}\|q_{\alpha,0}-q_{\beta,0}\|,
\qquad
t\ge 0.
\]
Consequently, every finite-$K$ diversity pattern created purely by initialization decays exponentially under the nonlinear flow.
\end{theorem}

\begin{proof}
Fix $\alpha,\beta$ and write $\Delta_t:=q_{\alpha,t}-q_{\beta,t}$. Then
\[
\dot\Delta_t
=
\mathcal C(\rho_t,q_{\alpha,t})-\mathcal C(\rho_t,q_{\beta,t}).
\]
Therefore
\[
\frac{d}{dt}\frac12\|\Delta_t\|^2
=
\langle \Delta_t,\dot\Delta_t\rangle
\le
-\gamma \|\Delta_t\|^2
\]
by \eqref{eq:one-sided-lipschitz}. Gronwall's inequality yields the claim.
\end{proof}

\begin{corollary}[Why the user's objection is mathematically correct]
\label{cor:user-correct}
If the long-time member dynamics are contractive in the sense of Theorem~\ref{thm:nonlinear-contraction}, then no one-shot initialization can guarantee persistent asymptotic disagreement. Sustained diversity requires at least one of:
\begin{itemize}
\item non-contractive member dynamics,
\item multiple attracting branches,
\item persistent centered forcing/noise,
\item explicit projection or regularization that preserves spread during training.
\end{itemize}
\end{corollary}

\section{What Can Be Guaranteed at a Converged Law}

Suppose now that the collapsed system has converged to a stationary point
\[
(\rho_\star,q_\star).
\]
Let the local predictor map at this stationary law be linearized in the member-specific state. Define the stationary Gramian
\begin{equation}
W_\star
:=
\int J_\star(x)^\top J_\star(x)\,\mu(\mathrm dx),
\qquad
J_\star(x):=\partial_q H(x;\rho_\star,q_\star),
\label{eq:stationary-gramian}
\end{equation}
for a chosen calibration distribution $\mu$.

\begin{remark}
This is the same Gramian as in the local theory, but now evaluated at the \emph{converged collapsed law}. That is the correct object if the design goal concerns the final predictor rather than the early trajectory.
\end{remark}

\begin{theorem}[Stationary-law optimal finite-$K$ design]
\label{thm:stationary-law-opt}
Fix a target rank $r\le K-1$ and a per-member deviation budget $\tau>0$. At the stationary law $(\rho_\star,q_\star)$, among all local member perturbation codes that
\begin{itemize}
\item are centered: $\sum_{\alpha=1}^K c_\alpha=0$,
\item satisfy the stationary per-member budget:
\[
\|W_\star^{1/2}c_\alpha\|^2\le \tau,
\qquad
\alpha=1,\dots,K,
\]
\end{itemize}
the best achievable minimum covered uncertainty eigenvalue is
\[
\frac{\tau}{r},
\]
and when $K=r+1$ it is attained by a simplex code built from the top-$r$ eigenspace of $W_\star$.
\end{theorem}

\begin{proof}
This is exactly the finite-$K$ robust-coverage theorem from the local geometry note, applied at the stationary Gramian $W_\star$. The centered condition preserves the ensemble mean at first order, and the simplex construction saturates the per-member bound while equalizing coverage across the protected eigendirections.
\end{proof}

\begin{corollary}[What is actually justified]
\label{cor:what-justified}
The local theory justifies the following statement:

\begin{quote}
If one wants the \emph{final} finite-$K$ ensemble to have stable mean prediction and balanced uncertainty coverage, then the diversity code should be designed from the stationary collapsed law, not merely from the initialization.
\end{quote}
\end{corollary}

\section{Practical Consequence: Two-Stage Training}

Theorem~\ref{thm:stationary-law-opt} suggests a concrete procedure.

\subsection*{Stage A: train a collapsed baseline}

Train with all member-specific parameters tied or identically initialized so that the system follows the collapsed baseline as closely as possible. Let the resulting terminal state be approximately
\[
(\hat\rho_\star,\hat q_\star).
\]

\subsection*{Stage B: estimate the stationary Gramian}

Using a calibration set, estimate
\[
\hat W_\star
=
\frac1{|B|}\sum_{x\in B}\hat J_\star(x)^\top \hat J_\star(x),
\qquad
\hat J_\star(x)
=
\partial_q H(x;\hat\rho_\star,\hat q_\star).
\]

\subsection*{Stage C: initialize a centered simplex code}

Let
\[
\hat W_\star e_i = \hat\lambda_i e_i,
\qquad
i=1,\dots,r.
\]
Choose simplex vertices $v_\alpha\in\R^r$ and set
\[
c_\alpha
=
\sqrt{\tau}\sum_{i=1}^r \hat\lambda_i^{-1/2}(v_\alpha)_i e_i.
\]
Then initialize
\[
q_{\alpha,0}^{\mathrm{stage\,C}}
=
\hat q_\star + c_\alpha.
\]

\subsection*{Stage D: constrained fine-tuning}

Fine-tune while enforcing
\[
\sum_{\alpha=1}^K c_{\alpha,t}=0,
\qquad
\| \hat W_\star^{1/2} c_{\alpha,t}\|^2 \le \tau
\]
either by projection or by reparameterization.

\begin{remark}[Why this is better than one-shot initialization]
This procedure targets the terminal law directly. The initial simplex code is no longer justified by ``it looked good at time zero'', but by ``it is locally optimal around the converged collapsed law we actually care about''.
\end{remark}

\section{What Remains Unproved}

The following are still not proved by the current notes:
\begin{itemize}
\item that a single early-time Gramian remains relevant throughout long nonlinear training;
\item that unconstrained fine-tuning preserves the stationary-law optimum;
\item that the estimated $\hat W_\star$ from finite data yields near-optimal performance;
\item that the final nonlinear fine-tuned solution remains in the local regime where the stationary-law linearization is accurate.
\end{itemize}

\begin{remark}
So the strongest defensible claim is not ``our initialization stays optimal forever'', but rather:

\begin{quote}
initialization-only design is generally transient; if one wants a principled final finite-$K$ ensemble, one should either design diversity from the converged collapsed law or actively maintain the desired geometry throughout training.
\end{quote}
\end{remark}

\section{Takeaway}

The user's concern is mathematically correct. Initialization gives only a local entrance condition. Long nonlinear training can erase it. To make a finite-$K$ design claim about the final predictor, the relevant object is the \emph{stationary law} and its Gramian, together with a training scheme that preserves zero-mean and bounded member deviation.

\end{document}
